% SOURCE docs/release/2026-10-29/source/docs/writing-repair/2026-10-29/src/72275a22fb653cf2.tex CHAR 0
\chapter{Continuous temporal uncertainty and an informative intervention}
\label{exec13:temporal-region}

\hypertarget{the-banking-question-and-the-observation-law}{%
\section{The banking question and the observation
law}\label{the-banking-question-and-the-observation-law}}

A bank can observe a stable association between an offer and a
subsequent outcome yet remain unable to price the effect of changing the
offer. The missing information can persist even with an exactly known
observational distribution. This chapter develops a restricted temporal
confidence construction and shows what an additional direct intervention
would contribute. The mathematical construction and exact calculations
are project derivations, not a fitted bank result.
\citep{exec13_temporal_argument,exec13_temporal_exact}

In the preserved temporal graph study, all valid cells covered and
retained their declared true graph in 512 of 512 requests. Weak
relevance and unrestricted latent observation each produced 512
whole-line sets. The persistent cell produced 507 finite sets and five
whole-line sets. Under false exclusion, the true graph was retained zero
times despite numerical effect coverage. These findings distinguish
numerical coverage from informative effects and from a valid causal
restriction. \citep{exec13_temporal_predecessor}

Let \(X_t=(Z_t,A_t,Y_t)'\) satisfy \[
 X_t=\rho X_{t-1}+\varepsilon_t,\quad
 \varepsilon_t\mathrel{\mathop\sim^{\rm iid}}N_3(0,\Sigma),\quad
 \Sigma\succ0,\quad \rho\in I\subset(-1,1).
\] Here \(I\) is a declared closed interval, not an estimated selected
lag. Initial observations have the invariant distribution
\(N_3(0,\Sigma/(1-\rho^2))\), independently of future innovations. We
observe \(X_0,\ldots,X_n\), with \(n\ge2\). The scalar lag common to all
coordinates and the Gaussian independent innovations are substantive
restrictions. A general lag matrix, unknown lag order, nonstationarity,
hidden temporal state, selection and measurement errors are outside this
law.

A structural innovation model is \[
 \varepsilon_Z=E_Z,\qquad
 \varepsilon_A=k\varepsilon_Z+U,\qquad
 \varepsilon_Y=\tau\varepsilon_A+c\varepsilon_Z+V,
\] where \(E_Z\) is independent of the jointly Gaussian pair \((U,V)\).
Dependence between \(U\) and \(V\) represents latent action/outcome
confounding. The target \(\tau\) is the contemporaneous response to
changing \(A_t\) while holding the past fixed. In observed-coordinate
structural equations, the lag matrix is \(\rho(I-B)\), where \(B\)
contains the contemporaneous coefficients. It is not permissible to
change that lag matrix to \(\rho I\) in the structural equations and
claim the same target/law.

The declared graph family may include the exclusion model \(c=0\) and a
broad latent model with unrestricted \(c\). In the former,
\(\Sigma_{ZY}=\tau\Sigma_{ZA}\). In the latter that equality need not
hold. An exclusion is a scientific restriction, not a conclusion from
fitting a favorable covariance. The construction below accommodates a
finite declared family through a union over graph/parameter pairs; it
does not claim an exhaustive search over unrestricted causal structures.

\hypertarget{a-feasible-region-with-continuous-unknown-lag}{%
\section{A feasible region with continuous unknown
lag}\label{a-feasible-region-with-continuous-unknown-lag}}

For every candidate \(r\in I\), form residuals \(E_t(r)=X_t-rX_{t-1}\).
Use the nine fixed directions \[
 \mathcal V=\{e_1,e_2,e_3,e_1+e_2,e_1-e_2,
 e_1+e_3,e_1-e_3,e_2+e_3,e_2-e_3\},\quad K=9.
\] For each direction, let \(S_v(r)\) be the ordinary sample variance of
\(v'E_t(r)\), with divisor \(\nu=n-1\). At the true lag only,
\(\nu S_v(\rho)/(v'\Sigma v)\) is the sum of squares of \(\nu\)
independent standard normal variables. To see this, scale the vector of
independent normal residuals by its population standard deviation and
apply an orthogonal matrix whose first row is the normalized all-ones
vector. The remaining \(\nu\) coordinates are independent standard
normals, and their squared norm is the centered sum of squares.
Consequently \[
 E[S_v(\rho)]=v'\Sigma v,\qquad
 \operatorname{Var}(S_v(\rho))=2(v'\Sigma v)^2/\nu.
\] For a declared rational tolerance \(0<t<1\), Markov's inequality
applied to the squared relative error gives \[
 P\{|S_v(\rho)/(v'\Sigma v)-1|>t\}\le\frac{2}{\nu t^2}.
\] If \(2K/(\nu t^2)\le\alpha_O\), the union bound supplies simultaneous
coverage at least \(1-\alpha_O\) for \[
 L_v(r)=\frac{S_v(r)}{1+t}\le v'\Sigma v
 \le \frac{S_v(r)}{1-t}=U_v(r),\qquad v\in\mathcal V,
\] evaluated at \(r=\rho\). Dependence between the nine directions is
harmless for this bound. Residuals at an incorrect lag need not be
independent.

Define \(\mathcal C_O\) to contain every tuple \((r,G,\lambda,\Sigma)\)
for which \(r\in I\), \(G\) is in the declared family, its parameters
\(\lambda\) generate \(\Sigma\succ0\) under the stated innovation law,
and all nine inequalities hold. All nuisance coefficients and the
continuous lag remain in the region. Let \[
 \mathcal S_O=\{\tau(G,\lambda):(r,G,\lambda,\Sigma)\in\mathcal C_O\}.
\] On simultaneous variance inclusion, the true tuple belongs to
\(\mathcal C_O\); hence its effect belongs to \(\mathcal S_O\). Thus
\(P(\tau_0\in\mathcal S_O)\ge1-\alpha_O\), provided the declared family
contains the true structural model. No multiple-testing penalty over the
continuum of candidate lags is required: this proof uses a single
coverage event at the fixed true lag, and the union includes that lag.
It does not declare every wrong-lag confidence calculation valid or
justify discarding the true lag by a separate selector.

The construction is deliberately conservative. With \(t=1/4\), \(K=9\)
and \(\alpha_O=.02\), it needs \(\nu\ge14{,}400\). The predecessor used
240 transitions; this particular finite bound cannot give it the
displayed two-sided relative band. At smaller sample size this profile
returns the whole line and a bound-insufficient disposition, rather than
inserting an invalid finite interval. The exact chi-square pivot above
offers a sharper candidate using suitably certified quantiles, but its
numerical implementation would need separate validation. No improvement
in cost or useful width is claimed from replacing a lag grid with this
conservative continuum argument.

\hypertarget{numerical-outer-bounds-that-cannot-lose-a-lag-or-a-component}{%
\section{Numerical outer bounds that cannot lose a lag or a
component}\label{numerical-outer-bounds-that-cannot-lose-a-lag-or-a-component}}

Let \(u_t=v'X_t\) and \(w_t=v'X_{t-1}\). Center each series over the
same \(n\) transitions. Direct expansion gives \[
 \begin{aligned}
 S_v(r)&=a_vr^2+b_vr+d_v,\\
 a_v&=\nu^{-1}\sum_t(w_t-\bar w)^2,\\
 b_v&=-2\nu^{-1}\sum_t(u_t-\bar u)(w_t-\bar w),\\
 d_v&=\nu^{-1}\sum_t(u_t-\bar u)^2.
 \end{aligned}
\] For a closed lag cell \(J=[r_L,r_U]\), its maximum is at an endpoint
because \(a_v\ge0\); its minimum is at an endpoint or at \(-b_v/(2a_v)\)
if that vertex lies in the cell. The constant case is handled
separately. Rational coefficients and cell endpoints therefore give
exact bounds \([s^-_v(J),s^+_v(J)]\). A finite partition covering all of
\(I\), with shared endpoints retained, is a continuum enclosure, not an
evaluation grid. Refining cells can reduce excess width; stopping early
preserves validity.

For exact supplied numbers, a conservative directional variance interval
over the cell is \([s^-_v/(1+t),s^+_v/(1-t)]\). Polarization gives an
entrywise covariance box. For \(i\ne j\), use \[
 \Sigma_{ij}\in\frac14
 [L_{e_i+e_j}(J)-U_{e_i-e_j}(J),
  U_{e_i+e_j}(J)-L_{e_i-e_j}(J)].
\] Dropping positive-definiteness and shared-parameter consistency
enlarges this box. It can make results much wider, but cannot remove the
true covariance. Positive-definiteness or graph constraints may narrow a
later implementation only through certified feasibility exclusion, not
by failure of an optimizer.

Exact arithmetic on decimal data alone does not enclose unknown
rounding. Suppose every coordinate of the ideal input is enclosed within
absolute error \(\eta\) of the supplied value, and let
\(h=\|v\|_1\eta(1+\max_{r\in J}|r|)\). Residual error is at most \(h\)
and its centered error at most \(2h\). Cauchy--Schwarz then bounds the
change in sample variance by \[
 B_v(J)=4h\sqrt{n/\nu}\sqrt{s^+_v(J)}+4nh^2/\nu.
\] Use rational upward bounds for the square roots and replace the
variance enclosure by \([\max(0,s^-_v-B_v),s^+_v+B_v]\) before dividing
by \(1\pm t\). This is an input-conversion certificate conditional on a
valid input enclosure. Undocumented measurement error, noisy proxies or
incorrect entity matches are changes to the observation law, not
floating-point errors that this formula repairs.

In an exclusion cell, write the resulting intervals as \(N_J\) for
\(\Sigma_{ZY}\) and \(D_J\) for \(\Sigma_{ZA}\). Retain the full set \[
 \mathcal T_J=\{\theta:\exists (N,D)\in N_J\times D_J,
                      \ N-\theta D=0\}.
\] When \(0\notin D_J\), its endpoints are the extreme corner ratios.
When the denominator straddles zero and the numerator excludes zero, two
rays can result. For \(N_J=[1,2]\) and \(D_J=[-1,1]\), the set is
\(( -\infty,-1]\cup[1,\infty)\). When both intervals include zero, their
rectangle contains \((0,0)\), so every \(\theta\) is feasible. When
\(D_J=\{0\}\) and \(0\notin N_J\), the moment set is empty. That is a
constraint inconsistency, not evidence for a zero effect.

Union these components over every lag cell and every retained graph. The
broad latent graph below contributes the whole line whenever its
observational law is admitted. Keep disconnected components unless they
overlap; never replace a pair of rays by a finite interval. The exact
reference checks ratio sets against the independent criterion that
\(N_J\) intersects \(\theta D_J\), including zero denominators and exact
endpoints. Bounds include every continuous lag by the quadratic proof;
sampled checks alone are not the certificate. The reference is not a
general graph feasibility solver and supplies no bank-data result.

\hypertarget{an-equal-law-witness-and-what-an-exclusion-actually-buys}{%
\section{An equal-law witness and what an exclusion actually
buys}\label{an-equal-law-witness-and-what-an-exclusion-actually-buys}}

Let \(Z,U,E\) be independent standard Gaussian innovations and set \[
 \varepsilon_Z=Z,\quad \varepsilon_A=Z+U,\quad
 \varepsilon_Y=\tau\varepsilon_A+(2-\tau)Z+(1-\tau)U+E.
\] For every real \(\tau\), the observed innovation is exactly
\((Z,Z+U,2Z+U+E)'\), with covariance \[
 \Sigma=\begin{pmatrix}1&1&2\\1&2&3\\2&3&6\end{pmatrix},
 \qquad\det\Sigma=1.
\] The confounded disturbance has variance \(1+(1-\tau)^2\) and
covariance \(1-\tau\) with \(U\); its conditional noise variance is
always one. Every world is a coherent positive-definite Gaussian law.
Attaching the same scalar lag and invariant initialization makes the
entire observed temporal law identical for every \(\tau\), while a
contemporaneous intervention on \(A\) changes \(Y\) at slope \(\tau\).
The equality is algebraic for all real effects, not merely an equality
across the finite checked examples.

The restriction \(c=0\) reduces this family to \(\tau=2\); strong
observed relevance \(\Sigma_{ZA}=1\) then gives the familiar ratio. But
relevance alone does not establish exclusion: all other worlds have the
same strong observed correlation with \(Z\). Merely observing more of
this law, using a different discovery optimizer or asking an LLM to
orient an edge cannot choose among these effects. No theorem here rules
out all additional measurements; it states exactly what the declared
observational law omits.

At zero relevance, even exclusion can leave every effect feasible. Take
\(\varepsilon_A=U\) and
\(\varepsilon_Y=\tau\varepsilon_A+(1-\tau)U+E=U+E\), with \(Z\)
independent. Now \(c=0\), \(\Sigma_{ZA}=\Sigma_{ZY}=0\), and the effect
is still arbitrary. A weak-instrument ratio that quietly deletes zero
cannot solve this information problem.

\hypertarget{a-direct-intervention-that-separates-effects-without-identifying-every-edge}{%
\section{A direct intervention that separates effects without
identifying every
edge}\label{a-direct-intervention-that-separates-effects-without-identifying-every-edge}}

Consider independent experimental episodes at a common fixed pre-action
history. Replace the action equation by \(A=a_0+\delta R\), where
\(\delta>0\) is known and \(R\in\{-1,1\}\) is assigned with \(m\)
episodes per arm independently of innovation roots. Keep the other
structural equations unchanged. Under the restricted linear Gaussian
model and an invariant \(\tau\), the outcome has the form \[
 Y=\mu+\tau\delta R+\xi,\quad
 \xi\mathrel{\mathop\sim^{\rm iid}}N(0,\sigma_I^2),\qquad \sigma_I^2>0.
\] The unknown \(\mu\) absorbs the fixed history and baseline action;
the unknown variance absorbs the remaining contemporaneous causes. The
episodes must be independent; taking adjacent dates from one continuing
customer trajectory does not satisfy this experiment. Identical variance
in the two arms and a common transported effect are additional model
assumptions.

Let \(\widehat\tau_I=(\bar Y_+-\bar Y_-)/(2\delta)\), and let \(S_p^2\)
be the usual pooled within-arm variance with \(\nu_I=2m-2\) degrees of
freedom. Orthogonal Gaussian projections within each arm make its sample
mean independent of its centered residuals. Hence \[
 T=\frac{\widehat\tau_I-\tau}{S_p/(\delta\sqrt{2m})}
   =\frac{Z_0}{\sqrt{Q/\nu_I}},\qquad
 Z_0\sim N(0,1),\quad Q\sim\chi^2_{\nu_I},\quad Z_0\perp Q.
\] For \(m>2\), integrating the chi-square density gives
\(E[Q^{-1}]=1/(\nu_I-2)\), so \(E[T^2]=\nu_I/(\nu_I-2)\). Choose any
rational \(b^2\ge\nu_I/((\nu_I-2)\alpha_I)\). Markov's inequality proves
finite coverage at least \(1-\alpha_I\) for \[
 \mathcal S_I=\left[\widehat\tau_I-
  \sqrt{\frac{b^2S_p^2}{2m\delta^2}},\quad
  \widehat\tau_I+\sqrt{\frac{b^2S_p^2}{2m\delta^2}}\right].
\] This conservative alternative avoids unverified numerical tail
quantiles. A sharper Student pivot interval is possible with certified
quantiles. Nonpositive pooled variance or insufficient degrees of
freedom produces a refusal/whole-line result under the declared
algorithm, never zero-width certainty. A nonnormal outcome model needs
another uncertainty argument. If input numbers are enclosed rather than
exact, enclose the arm means and pooled variance as well before
projecting these endpoints.

Under a common effect parameter, report
\(\mathcal S=\mathcal S_O\cap\mathcal S_I\). The two component coverage
events need not be independent: the union bound gives at least
\(1-\alpha_O-\alpha_I\). Using outer numerical supersets preserves this
inclusion proof. An empty intersection is a retained
inconsistency/failure event, not an invitation to change the error
budget or delete a graph after seeing the result. If no observational
restriction is defensible, use \(\mathcal S_I\) alone under its own
experiment and target assumptions.

In the equal-law family, intervention changes the arm mean difference to
\(2\delta\tau\). Distinct effects are now distinguishable, although
several graphs can remain compatible. Identifying the decision-relevant
effect does not require identifying every edge. Randomizing a variable
that merely encourages the action is a different experiment: exclusion,
compliance and relevance then need their own argument.

\hypertarget{precision-costs-and-the-exact-remaining-scope}{%
\section{Precision, costs and the exact remaining
scope}\label{precision-costs-and-the-exact-remaining-scope}}

For a binary offer, take \(a_0=\delta=1/2\). A planning illustration
with \(m=200\) per arm, \(\alpha_I=.03\), a reported effect of twenty
dollars and \(S_p^2=100\) dollars squared gives \(b^2=9950/297\) and
half-width \(\sqrt{9950/297}\), approximately 5.788 dollars. The exact
reference uses a rational upward square-root enclosure. These are
supplied planning summaries, not observations from a completed trial.

For 3,000 eligible customers and fifteen dollars incremental
implementation cost per offer, the point scenario has fifteen thousand
dollars net benefit. Its conservative interval still includes a small
loss. Scaling an effect by population does not remove uncertainty, and
changing the intervention or customer population may invalidate
transport. A genuine next-experiment proposal should ask whether
resolving this decision is worth the experiment's cost, compare a
cheaper operational alternative, and include delayed outcomes,
noncompliance and treatment costs in the estimand. No measured bank ROI
is inferred from this illustration.

What is resolved is the mathematical route for one continuous scalar-lag
Gaussian family, a certified enclosure strategy, exact information
counterexamples and a separating intervention with an explicit finite
bound. What remains is a resource-profiled implementation and
prospectively registered study, tighter useful uncertainty, justified
input and model assumptions, actual interventions/measurements if the
proposal advances, and the general
unknown-lag-matrix/latent-state/non-Gaussian/heterogeneous-effect
problem. This argument establishes no general temporal-graph capability,
bank validation or independent acceptance. The computational and
information routes are concrete enough to design the next bounded
public/generated experiment. The following design analysis tests whether the displayed bounds could be informative at the predecessor's sample size and whether its qualification screen has adequate power. It keeps theorem validity, finite availability and a useful financial comparison as separate requirements.

\label{exec13:temporal-region-end}


% SOURCE docs/release/2026-10-29/source/docs/writing-repair/2026-10-29/src/db9ea6894074d6d4.tex CHAR 0
\section{When a valid temporal region is still an unsuitable study design}
\label{exec13:temporal-design-review}

The restricted continuous-lag argument separates mathematical coverage
from decision-useful precision. At the true common scalar lag, invariant
Gaussian initialization and independent innovations justify the projected
variance argument. A union over candidate lags need only retain the true
lag on that event; incorrect lags need not themselves produce independent
residuals. This is a useful restricted result. It does not establish a
confidence procedure for an arbitrary lag matrix, unknown order,
nonstationary hidden histories or a general learned partial ancestral
graph. These differences also determine whether a proposed numerical
study is worth conducting \citep{exec13_temporal_design}.

Consider the displayed Chebyshev construction with nine variance
directions, relative tolerance $t=1/4$ and observational error allocation
$\alpha_O=.02$. For a Gaussian sample variance with $\nu$ residual
degrees of freedom, the relative variance is $2/\nu$. Chebyshev and a
union bound therefore give
\[
 \Pr(\hbox{at least one relative deviation exceeds }t)
 \leq \frac{18}{\nu t^2}.
\]
To make this bound at most $.02$ requires $\nu\geq14{,}400$, or
$14{,}401$ transitions under the stated residual convention. With $240$
transitions and $239$ degrees of freedom, the untruncated expression is
approximately $1.20502$ and gives no useful coverage bound. The specified
whole-line fallback is honest about this limitation. Increasing the number
of Monte Carlo replications cannot make that per-request fallback
decision-useful \citep{exec13_temporal_design}.

A sharper pivot is a concrete development route. At the restricted true
lag, an exact chi-square variance pivot can replace the Chebyshev band
while preserving the continuous union. Ordinary development quantiles
at $239$ degrees of freedom and the nine-way $.02$ allocation give
variance multipliers approximately $[.76731,1.34533]$. These values
indicate potential improvement; they are not certified release endpoints.
The implementation must enclose the tail quantiles, observable inputs
and propagated covariance constraints, preserve conservative fallback
under timeout, and profile the continuous-lag calculation. Sampling a
lag grid or accepting a point optimizer's failure does not certify the
whole continuum \citep{exec13_temporal_design}.

The preserved finite-lag study also requires a precise reading. Its
persistent cell returned $507$ finite sets among $512$ requests, so
$507/512=.990234375$ passes the original raw finite-fraction threshold
of $.99$. The five whole-line answers remain visible, but do not justify
declaring that particular threshold failed. Weak relevance and the
unrestricted latent observational cell each returned zero finite sets
and still fail it. False exclusion retained the true graph in zero
requests, despite wide effect sets covering the numerical effect.
These findings concern different obligations: availability, effect
coverage, and scientific validity of the graph class
\citep{exec13_temporal_design}.

A lower confidence-bound rule for availability would be a different
criterion. At family error $.05$ across sixteen comparisons, the
one-sided lower bound for $507/512$ is about $.97130$; even $512/512$
gives about $.98880$. The historical study used a lower-bound screen
for coverage, but a raw fraction for finite availability. They should
not be interchanged retrospectively. A prospective successor may
justify different criteria, provided it states them before final data
and keeps the predecessor's outcomes under their original definitions
\citep{exec13_temporal_design}.

The coverage screen has its own design-power issue. With $512$
independent final requests per cell and error $.05/16$, at least $500$
successes are required for the exact one-sided lower bound to exceed
or equal $.95$. Exact binomial-tail calculations give per-cell passing
probabilities approximately $.001849$, $.234929$, $.770054$ and
$.997670$ when true coverage is respectively $.95$, $.97$, $.98$ and
$.99$. Thus a nominal $.95$ method is very unlikely to certify a
lower bound of $.95$ using this screen. More replications improve
discrimination above the boundary; they do not create high power at
the boundary while retaining the same error control. No product of
these per-cell powers is a joint power calculation without an
across-cell dependence argument \citep{exec13_temporal_design}.

Useful precision also requires an information distinction. In the
unrestricted equal-law latent family, every real effect can produce
the same observed temporal law. No numerical improvement can create
a uniformly finite informative effect set from those unchanged
observations. A direct assigned-action experiment changes the
information available. The proposed Gaussian episode argument assumes
independent episodes, a common transported effect and equal arm
variances. Adjacent records from one trajectory or encouragement in
place of the action require different assumptions and derivations
\citep{exec13_temporal_design}.

The supplied-summary example makes the planning distinction concrete.
At $200$ episodes per arm, reported effect $20$ dollars, reported
pooled variance $100$ dollars squared and the specified conservative
bound, the half-width is about $5.78806$ dollars. For $3{,}000$
eligible customers and an incremental cost of $15$ dollars per offer,
the net-benefit lower endpoint is approximately $-2{,}364.19$ dollars.
Holding the same reported variance and effect fixed, $268$ episodes
per arm first reduce the half-width below $5$ dollars, while $400$
give about $4.08761$. In a future experiment both the estimated effect
and pooled variance would vary. These are conditional arithmetic
screens, not prospective power or recommended bank sample sizes
\citep{exec13_temporal_design}.

The next restricted study therefore needs a certified executable and
input-enclosure policy, measured resource limits, a fixed final grid
of observational and added-information regimes, and an explicit
useful-width or loss criterion. A generated experiment can develop
these elements now. It should preserve strong, weak and zero relevance,
false exclusion, latent equal-law cases and all interval topologies.
Final seeds, replication budgets and stopping rules must follow the
design review. General hidden-state, non-Gaussian, selection and
heterogeneous-transport questions remain unresolved beyond the
restricted claim; public observations alone cannot decide among
observationally equivalent causal worlds
\citep{exec13_temporal_design}. The next chapter adds fitted prediction to the information budget. Its pilot certificates can control specified nuisance errors, but the same questions about weak identification, dependence and useful width still determine whether a larger experiment is justified.
\label{exec13:temporal-design-review-end}

